\chapter{Planificación del proyecto}
\label{chap:planificacion}

\ph{Este capítulo describe cómo se ha gestionado el proyecto: metodología de trabajo,
fases, hitos, riesgos y presupuesto. Demuestra que el TFG no es solo un producto
software, sino un proyecto planificado y controlado.}

\section{Metodología de desarrollo}
\label{sec:metodologia}

\ph{Justifica la metodología (o metodologías) elegida(s) para el desarrollo y para
la redacción de la memoria. Es habitual usar enfoques distintos para cada fase.}

\ph{Puntos a cubrir:}
\begin{itemize}
  \item \ph{Metodología elegida (Waterfall, Scrum, Kanban, iterativa\dots) y por qué.}
  \item \ph{Cómo se gestiona el control de versiones (Git, GitHub/GitLab\dots).}
  \item \ph{Si se usan herramientas de seguimiento (issues, tableros, CI), nómbralas brevemente.}
\end{itemize}

\section{Fases e hitos}
\label{sec:fases}

\ph{Divide el proyecto en fases claras con un hito (entregable verificable) por fase.
Es muy recomendable comparar la estimación inicial con las horas realmente invertidas,
y comentar las desviaciones.}

\ph{Tabla de fases del desarrollo: rellena las filas con tus propias fases.}

\begin{longtable}{@{}P{0.4cm} P{4cm} P{2cm} P{2cm} P{6cm}@{}}
  \toprule
  \textbf{\#} & \textbf{Fase} & \textbf{Est.} & \textbf{Real} & \textbf{Hito / Entregable} \\
  \midrule
  \endhead
  1 & \ph{Fase 1} & \ph{X h} & \ph{Y h} & \ph{Entregable de la fase 1.} \\[0.3em]
  2 & \ph{Fase 2} & \ph{X h} & \ph{Y h} & \ph{Entregable de la fase 2.} \\[0.3em]
  3 & \ph{Fase 3} & \ph{X h} & \ph{Y h} & \ph{Entregable de la fase 3.} \\[0.3em]
  4 & \ph{Fase 4} & \ph{X h} & \ph{Y h} & \ph{Entregable de la fase 4.} \\[0.3em]
  \midrule
  & \textbf{Total} & \textbf{\ph{X h}} & \textbf{\ph{Y h}} & \\
  \bottomrule
  \caption{Fases del proyecto e hitos asociados (estimadas vs.\ reales).}
  \label{tab:fases-desarrollo}
\end{longtable}

\ph{Comenta la desviación: en qué fase se concentró, por qué se produjo y cómo se gestionó.}

\ph{Tabla equivalente para la redacción de la memoria.}

\begin{longtable}{@{}P{0.4cm} P{4cm} P{2cm} P{2cm} P{6cm}@{}}
  \toprule
  \textbf{\#} & \textbf{Capítulo} & \textbf{Est.} & \textbf{Real} & \textbf{Hito / Entregable} \\
  \midrule
  \endhead
  1 & Introducción              & \ph{X h} & \ph{Y h} & \ph{Capítulo redactado y revisado.} \\[0.3em]
  2 & Estado del arte           & \ph{X h} & \ph{Y h} & \ph{Capítulo redactado y revisado.} \\[0.3em]
  3 & Planificación             & \ph{X h} & \ph{Y h} & \ph{Capítulo redactado y revisado.} \\[0.3em]
  4 & Análisis de requisitos    & \ph{X h} & \ph{Y h} & \ph{Capítulo redactado y revisado.} \\[0.3em]
  5 & Diseño de la solución     & \ph{X h} & \ph{Y h} & \ph{Capítulo redactado y revisado.} \\[0.3em]
  6 & Implementación            & \ph{X h} & \ph{Y h} & \ph{Capítulo redactado y revisado.} \\[0.3em]
  7 & Pruebas y validación      & \ph{X h} & \ph{Y h} & \ph{Capítulo redactado y revisado.} \\[0.3em]
  8 & Conclusiones              & \ph{X h} & \ph{Y h} & \ph{Capítulo redactado y revisado.} \\
  \midrule
  & \textbf{Total} & \textbf{\ph{X h}} & \textbf{\ph{Y h}} & \\
  \bottomrule
  \caption{Fases de redacción de la memoria (estimadas vs.\ reales).}
  \label{tab:fases-memoria}
\end{longtable}

\ph{Tabla de resumen agrupada (planificación + memoria + desarrollo + revisión).}

\begin{longtable}{@{}P{0.4cm} P{4cm} P{2cm} P{2cm} P{6cm}@{}}
  \toprule
  \textbf{\#} & \textbf{Fase} & \textbf{Est.} & \textbf{Real} & \textbf{Hito / Entregable} \\
  \midrule
  \endhead
  1 & Planificación & \ph{X h} & \ph{Y h} & \ph{Hito.} \\[0.3em]
  2 & Memoria       & \ph{X h} & \ph{Y h} & \ph{Hito.} \\[0.3em]
  3 & Desarrollo    & \ph{X h} & \ph{Y h} & \ph{Hito.} \\[0.3em]
  4 & Revisión      & \ph{X h} & \ph{Y h} & \ph{Hito.} \\[0.3em]
  \midrule
  & \textbf{Total} & \textbf{\ph{X h}} & \textbf{\ph{Y h}} & \\
  \bottomrule
  \caption{Resumen global del proyecto (estimadas vs.\ reales).}
  \label{tab:fases-conjunta}
\end{longtable}

\ph{Cierra la sección con una valoración global de la desviación temporal.}

\section{Gestión de riesgos}
\label{sec:riesgos}

\ph{Identifica los riesgos principales que podían comprometer el cumplimiento de los
objetivos o de los plazos. Para cada uno indica la probabilidad estimada,
el impacto y el plan de mitigación.}

\begin{longtable}{P{4.5cm} P{2.2cm} P{2.2cm} P{5cm}}
  \toprule
  \textbf{Riesgo} & \textbf{Probabilidad} & \textbf{Impacto} & \textbf{Mitigación} \\
  \midrule
  \endhead
  \ph{Riesgo 1} & \ph{Alta} & \ph{Alto} & \ph{Plan de mitigación.} \\[0.3em]
  \ph{Riesgo 2} & \ph{Media} & \ph{Medio} & \ph{Plan de mitigación.} \\[0.3em]
  \ph{Riesgo 3} & \ph{Baja} & \ph{Alto} & \ph{Plan de mitigación.} \\[0.3em]
  \ph{Riesgo 4} & \ph{Media} & \ph{Bajo} & \ph{Plan de mitigación.} \\
  \bottomrule
  \caption{Riesgos identificados y planes de mitigación.}
  \label{tab:riesgos}
\end{longtable}

\section{Presupuesto}
\label{sec:presupuesto}

\ph{Presenta el presupuesto del proyecto. Una opción habitual es seguir el modelo
TCO (\textit{Total Cost of Ownership}), separando costes de capital (CapEx) y costes
operacionales (OpEx). Conviene incluir, junto al coste real (a menudo 0\,€ en contexto
académico), el valor de mercado equivalente para una valoración realista.}

\subsection*{CapEx --- Costes de capital}

\ph{Inversiones en infraestructura física y licencias de uso prolongado.}

\begin{longtable}{P{4cm} P{4cm} P{2.5cm} P{2.5cm}}
  \toprule
  \textbf{Recurso} & \textbf{Descripción} & \textbf{Coste real} & \textbf{Valor} \\
  \midrule
  \endhead
  \ph{Recurso 1} & \ph{Descripción.} & \ph{0,00\,€} & \ph{X\,€/año} \\[0.3em]
  \ph{Recurso 2} & \ph{Descripción.} & \ph{0,00\,€} & \ph{X\,€}     \\[0.3em]
  \midrule
  \textbf{Total CapEx} & & \textbf{\ph{0,00\,€}} & \textbf{\ph{X\,€}} \\
  \bottomrule
  \caption{Costes de capital (CapEx).}
  \label{tab:capex}
\end{longtable}

\ph{Justifica brevemente cómo se ha estimado el valor de mercado de cada recurso
(por ejemplo, tarifa pública de un proveedor cloud o precio comercial de una licencia).}

\subsection*{OpEx --- Costes operacionales}

\ph{Gastos recurrentes asociados al desarrollo y mantenimiento.}

\subsubsection*{Personal}

\ph{Indica el rol, la tarifa por hora y la justificación de esa tarifa
(convenio colectivo, mercado laboral de referencia, etc.).}

\begin{longtable}{P{3.5cm} P{2cm} P{2cm} P{3cm} P{3cm}}
  \toprule
  \textbf{Actividad} & \textbf{H. estimadas} & \textbf{H. reales} &
  \textbf{Coste estimado (€)} & \textbf{Coste real (€)} \\
  \midrule
  \endhead
  \ph{Actividad 1}        & \ph{X} & \ph{Y} & \ph{0,00} & \ph{0,00} \\[0.3em]
  \ph{Actividad 2}        & \ph{X} & \ph{Y} & \ph{0,00} & \ph{0,00} \\[0.3em]
  \ph{Actividad 3}        & \ph{X} & \ph{Y} & \ph{0,00} & \ph{0,00} \\[0.3em]
  \midrule
  \textbf{Total personal} & \textbf{\ph{X}} & \textbf{\ph{Y}} & \textbf{\ph{0,00}} & \textbf{\ph{0,00}} \\
  \bottomrule
  \caption{Desglose de costes de personal (OpEx).}
  \label{tab:coste_personal}
\end{longtable}

\subsubsection*{Herramientas y servicios de desarrollo}

\begin{longtable}{P{4cm} P{4cm} P{1.8cm} P{2.7cm}}
  \toprule
  \textbf{Herramienta} & \textbf{Licencia utilizada} & \textbf{Coste real} & \textbf{Valor de mercado} \\
  \midrule
  \endhead
  \ph{Herramienta 1} & \ph{Licencia.} & \ph{0\,€} & \ph{X\,€/año} \\[0.3em]
  \ph{Herramienta 2} & \ph{Licencia.} & \ph{0\,€} & \ph{X\,€/año} \\[0.3em]
  \midrule
  \textbf{Total herramientas} & & \textbf{\ph{0\,€}} & \textbf{\ph{X\,€/año}} \\
  \bottomrule
  \caption{Costes de herramientas y servicios de desarrollo (OpEx).}
  \label{tab:coste_herramientas}
\end{longtable}

\subsection*{Resumen del presupuesto (TCO)}

\begin{longtable}{P{4cm} P{3cm} P{3cm} P{3cm}}
  \toprule
  \textbf{Categoría} & \textbf{Tipo} & \textbf{Coste real} & \textbf{Valor de mercado} \\
  \midrule
  \endhead
  Infraestructura y licencias & CapEx & \ph{0,00\,€} & \ph{X\,€} \\
  Personal                    & OpEx  & \ph{X\,€}    & \ph{X\,€} \\
  Herramientas y servicios    & OpEx  & \ph{0,00\,€} & \ph{X\,€} \\
  \midrule
  \textbf{TCO total}          &       & \textbf{\ph{X\,€}} & \textbf{\ph{X\,€}} \\
  \bottomrule
  \caption{Resumen del presupuesto del proyecto (TCO = CapEx + OpEx).}
  \label{tab:presupuesto_total}
\end{longtable}

\ph{Cierra la sección comentando el coste real frente al valor de mercado
y la influencia del entorno académico (licencias educativas, equipos cedidos\dots).}
